% Source: PARADIGM.md section 6c (Proposition H.7, Remark H.7a, Theorem I, Proposition H.8, route to Conjecture H.6,
% evidence).  Notation: Section~\ref{sec:setting} and Section~\ref{sec:thmH}.
\section{Finite gain: semi-gradient and non-reciprocal structure}\label{sec:finitegain}

This section continues Section~\ref{sec:thmH} at finite feedback gain $k_e$. The notation is that of
Section~\ref{sec:setting} and Section~\ref{sec:thmH}: the saturated features $\tilde\Phi$, the residual $r_\mu$, the
potential $V_L[\mu]$, the averaged system $\Aavg$ and the law $(\mathrm{L}_\tau)$. We write
$\tilde\Phi_\vartheta := \tilde\Phi(\cdot;\vartheta)$.

\subsection{Frozen-measure structure of the exactly averaged loop}

\begin{thm}{Proposition}{H.7}[frozen-measure structure of the exactly averaged closed loop]
Assume the following:
\begin{itemize}
  \item $f\in C^2(\R^n)$ is autonomous;
  \item $x_m\in C^2$ is $T_0$-periodic, with oriented orbit $\mathcal{O}$;
  \item the features are the saturated features of the design (Section~\ref{sec:thmH});
  \item $\mu\in\Prob(\R^p)$ is frozen.
\end{itemize}
Then for every $R>0$ there is $k_0(R)$, independent of $\mu$, such that for $k_e\ge k_0(R)$ the following hold.
\begin{enumerate}[label=(\alph*)]
  \item \emph{Periodic error.} The frozen-$\mu$ error equation $\dot e = r_\mu(x_m+e) - k_e e$ has exactly one
    $T_0$-periodic solution $e_\mu$ with $\sup_t\norm{e_\mu}\le R$. It attracts exponentially every solution with
    $\norm{e(0)}\le R$, and
    \begin{equation}\label{eq:finitegain:emu}
      e_\mu = \frac{r_\mu(x_m)}{k_e}
        + \frac{\nabla r_\mu(x_m)\, r_\mu(x_m) - \frac{\mathrm{d}}{\mathrm{d}t} r_\mu(x_m(t))}{k_e^2}
        + O(k_e^{-3}),
    \end{equation}
    uniformly in $t$ and $\mu$.
    \begin{itemize}
      \item Here $\nabla r_\mu$ is the Jacobian $(\partial r_k/\partial x_l)$, and
        $\partial_x\tilde\Phi = (\partial\tilde\Phi_k/\partial x_l)$.
      \item If $f$ is one-sided Lipschitz, $\ip{f(x)-f(y)}{x-y}\le L_f\norm{x-y}^2$, then for
        $k_e > L_f + a_{\max}\sqrt{n}\,e^{-1/2}/s_{\min}$ the periodic solution is unique and attracts every
        solution.
      \item Without such a condition uniqueness is only local. For $f(x)=x^2$ there is a second, repelling periodic
        solution near $e\approx k_e$, and solutions above it blow up.
    \end{itemize}

  \item \emph{Expansion.} The period average of the drift of $(\mathrm{L}_\tau)$ along $e_\mu$ is
    $G\nabla_\vartheta P_\mu(\vartheta) - \sigma(\vartheta-\vartheta^0)$, where
    $P_\mu(\vartheta) := \langle e_\mu^\T\tilde\Phi(x_m+e_\mu;\vartheta)\rangle_t$ satisfies
    \begin{gather}
      P_\mu(\vartheta) = -\frac{1}{k_e}V_L[\mu](\vartheta)
        + \frac{1}{k_e^2}\bigl(U_{\mathrm{even}}[\mu](\vartheta) + U_{\mathrm{odd}}[\mu](\vartheta)\bigr)
        + O(k_e^{-3}),\label{eq:finitegain:Pmu}\\
      U_{\mathrm{even}}[\mu](\vartheta) := \int(r_\mu\cdot\nabla_x)\bigl(r_\mu^\T\tilde\Phi(\cdot;\vartheta)\bigr)\dd\rho,
      \qquad
      U_{\mathrm{odd}}[\mu](\vartheta) := -\frac{1}{T_0}\oint_{\mathcal{O}}\tilde\Phi(\cdot;\vartheta)^\T\dd r_\mu,
      \label{eq:finitegain:Ueo}
    \end{gather}
    uniformly in $\vartheta$ and $\mu$, together with $\vartheta$-gradients.

  \item \emph{Semi-gradient form.}
    \begin{equation}\label{eq:finitegain:semigrad}
      -\frac{1}{k_e}V_L[\mu] + \frac{1}{k_e^2}U_{\mathrm{even}}[\mu]
        = \frac{1}{k_e}\int r_\mu^\T\tilde\Phi(\cdot;\vartheta)\dd\rho^{\mathrm{cl}}_\mu + O(k_e^{-3}).
    \end{equation}
    Here $\rho^{\mathrm{cl}}_\mu$ is the occupation measure of the closed-loop state $x_m+e_\mu$.
    \begin{itemize}
      \item So, apart from the lag term $U_{\mathrm{odd}}$, the averaged law follows the gradient of the closed-loop
        loss $\frac12\int\norm{r_\nu}^2\dd\rho^{\mathrm{cl}}_\mu$ at $\nu=\mu$, with the closed-loop state
        distribution $\rho^{\mathrm{cl}}_\mu$ held fixed.
      \item This is a semi-gradient: the update freezes part of its own parameter dependence. In TD learning the
        frozen part is the bootstrapped target, and the data distribution does not depend on the
        parameters~\cite{S11,S10}.
      \item Here the frozen part is the data distribution $\rho^{\mathrm{cl}}_\mu$. This is the structure of
        performative prediction~\cite{S14}: steps taken with the distribution held at its current value converge to
        performatively stable points, not to optima.
      \item \textbf{The performative defect carries the even non-reciprocity.} With
        $J^{\mathrm{cl}}(\mu) := \frac12\int\norm{r_\mu}^2\dd\rho^{\mathrm{cl}}_\mu$,
        \begin{equation}\label{eq:finitegain:defect}
          \begin{aligned}
            &-\frac{1}{k_e}V_L[\mu] + \frac{1}{k_e^2}U_{\mathrm{even}}[\mu]\\
            &\qquad= -\frac{1}{k_e}\frac{\delta J^{\mathrm{cl}}}{\delta\mu}
              - \frac{1}{k_e^2}\int(\tilde\Phi_\vartheta\cdot\nabla)\bigl(\tfrac12\norm{r_\mu}^2\bigr)\dd\rho
              + O(k_e^{-3}).
          \end{aligned}
        \end{equation}
        The first term is a true first variation, hence reciprocal. The second is the defect: it has the same order
        as $U_{\mathrm{even}}$, and it carries all of $A_{\mathrm{even}}$ in~(e).
    \end{itemize}

  \item \emph{Circulation kernel.}
    $U_{\mathrm{odd}}[\mu](\vartheta) = -\frac{1}{T_0}\oint_{\mathcal{O}}\tilde\Phi(\cdot;\vartheta)^\T\dd f
      + \int K(\vartheta,\vartheta')\dd\mu(\vartheta')$, with
    \begin{equation}\label{eq:finitegain:K}
      K(\vartheta,\vartheta') := \frac{1}{T_0}\oint_{\mathcal{O}}\tilde\Phi(\cdot;\vartheta)^\T\dd\tilde\Phi(\cdot;\vartheta')
        = \Bigl\langle\tilde\Phi(x_m;\vartheta)^\T\,\frac{\mathrm{d}}{\mathrm{d}t}\tilde\Phi(x_m;\vartheta')\Bigr\rangle_t .
    \end{equation}
    \begin{itemize}
      \item $K$ is antisymmetric.
      \item By Stokes,
        $K(\vartheta,\vartheta') = \frac{1}{T_0}\int_\Sigma\sum_k\mathrm{d}\tilde\Phi_k(\cdot;\vartheta)\wedge
        \mathrm{d}\tilde\Phi_k(\cdot;\vartheta')$ for any $2$-chain $\Sigma$ with $\partial\Sigma=\mathcal{O}$.
      \item For $n=2$ and a simple, positively oriented orbit bounding $D$,
        \[
          K(\vartheta,\vartheta') = \frac{1}{T_0}\sum_k\int_D\det\bigl[\nabla_x\tilde\Phi_k(\cdot;\vartheta),
          \nabla_x\tilde\Phi_k(\cdot;\vartheta')\bigr]\dd A.
        \]
      \item $K$ is an integral over the enclosed region, not the area itself. For Gaussian features,
        $\sum_k\det[\nabla\tilde\Phi_k(\vartheta),\nabla\tilde\Phi_k(\vartheta')]
        = (a^\T a')\cdot\det[\nabla\varphi,\nabla\varphi']$, so $K=0$ whenever $a^\T a'=0$ or $c=c'$, whatever the
        area.
      \item Hence $U_{\mathrm{odd}}\equiv0$ for $n=1$.
      \item Reversing the orbit preserves $\rho$, flips $U_{\mathrm{odd}}$, and leaves $U_{\mathrm{even}}$ unchanged.
    \end{itemize}

  \item \emph{Non-reciprocity.} Write $S(\vartheta,\vartheta') := \delta P_\mu(\vartheta)/\delta\mu(\vartheta')$ for
    the derivative in the direction $\delta_{\vartheta'}$, and call $S(\vartheta,\vartheta') - S(\vartheta',\vartheta)$
    its antisymmetric part (without the factor $\frac12$). Adding a function of $\vartheta$ alone, or a constant
    $c(\mu)$, to $P_\mu$ changes the antisymmetric part by a coboundary, which the criterion below allows.
    \begin{itemize}
      \item At order $1/k_e$, $S$ is symmetric: it equals
        $-\frac{1}{k_e}\int\tilde\Phi(\cdot;\vartheta)^\T\tilde\Phi(\cdot;\vartheta')\dd\rho$.
      \item At order $1/k_e^2$ it is not. The second variation of the $k_e^{-2}$ coefficient
        $U_{\mathrm{even}}+U_{\mathrm{odd}}$ has antisymmetric part $A_{\mathrm{even}}+2K$, with
        \begin{equation}\label{eq:finitegain:Aeven}
          A_{\mathrm{even}}(\vartheta,\vartheta') := -\int\Bigl[(\tilde\Phi_{\vartheta'}\cdot\nabla)
            (r_\mu^\T\tilde\Phi_\vartheta) - (\tilde\Phi_\vartheta\cdot\nabla)(r_\mu^\T\tilde\Phi_{\vartheta'})\Bigr]
            \dd\rho .
        \end{equation}
      \item $P_\mu$, up to $O(k_e^{-3})$, can be a first variation of a functional of $\mu$ only if
        $A_{\mathrm{even}}+2K$ is a coboundary $a(\vartheta)-a(\vartheta')$, that is, only if all its cyclic sums
        vanish.
      \item They do not vanish in the checked instances, for $n=1$ (from $A_{\mathrm{even}}$ alone) and for $n=2$.
        Theorem I shows that the cyclic condition is equivalent to $A_{\mathrm{even}}+2K\equiv0$, and characterises
        when that happens: for $n=1$ exactly when the residual vanishes on the orbit; for $n\ge2$ at perfect fit
        exactly when the orbit carries no circulation and the residual's Jacobian is symmetric along the orbit. It
        also shows (Theorem I(vi)) that no open set of measures satisfies the condition, for any $n$.
      \item So at finite $k_e$, the exactly averaged loop is a non-reciprocal McKean--Vlasov dynamics: it is not the
        $\WG$-gradient flow of any free energy on any open set of measures. A different, $\mu$-dependent metric is
        not excluded by this argument.
    \end{itemize}
\end{enumerate}
\end{thm}
\status{The frozen-measure part is proved. Non-reciprocity of the exact second variation is checked to 3\%
(asserted; 1--1.5\% measured) and characterised by Theorem I. The Stokes identity and the parity under orbit
reversal hold to machine precision. The expansion is checked through its $O(1/k_e)$ convergence. Novelty (literature audit): no precedent found
(targeted search); ingredients known (higher-order averaging; performative prediction~\cite{S14}; semi-gradient
TD~\cite{S11,S10}).}

\begin{proof}
\emph{(a)} Write $e = r_\mu(x_m)/k_e + w$. Then $w$ solves $\dot w = -k_e w + g(t,w)$ with
$g = \nabla r_\mu(x_m)e - \dot r_\mu/k_e + O(\abs{e}^2)$.
\begin{itemize}
  \item The tube $\{\norm{e}\le R\}$ is forward invariant once $k_eR > \sup_{\norm{e}\le R}\norm{r_\mu(x_m+e)}$.
  \item Inside the tube, solutions contract at rate $k_e - \Lip_R(r_\mu)$, which gives existence, uniqueness and
    attraction there. A contraction on $T_0$-periodic functions gives the expansion.
  \item The bounds on $r_\mu$, $\nabla r_\mu$ and $\nabla^2 r_\mu$ on the $R$-neighbourhood of the orbit are uniform in
    $\mu$, by the saturation and $f\in C^2$.
  \item The global statement uses $L_f$ in place of the local Lipschitz constant of $f$.
\end{itemize}

\noindent\emph{(b)} Taylor-expand
$e^\T\tilde\Phi(x_m+e;\vartheta) = e^\T\tilde\Phi(x_m;\vartheta) + e^\T\partial_x\tilde\Phi(x_m;\vartheta)e
+ O(\abs{e}^3)$, and insert~(a).
\begin{itemize}
  \item The period averages of the time-symmetric terms are integrals against $\rho$.
  \item $(\nabla r\, r)^\T\tilde\Phi + r^\T(\partial_x\tilde\Phi)r = (r\cdot\nabla)(r^\T\tilde\Phi)$.
  \item $\vartheta$-derivatives commute with the average.
  \item The remainders are uniform, because the $x$- and $\vartheta$-derivatives of $\tilde\Phi$ of the orders needed
    are bounded.
\end{itemize}

\noindent\emph{(c)} We have
\[
  \int r_\mu^\T\tilde\Phi_\vartheta\dd\rho^{\mathrm{cl}}_\mu
  = \bigl\langle r_\mu(x_m+e_\mu)^\T\tilde\Phi(x_m+e_\mu;\vartheta)\bigr\rangle
  = \int r_\mu^\T\tilde\Phi_\vartheta\dd\rho
    + \bigl\langle\nabla_x(r_\mu^\T\tilde\Phi_\vartheta)\cdot e_\mu\bigr\rangle + O(k_e^{-2}).
\]
By~(a) this equals $\int r_\mu^\T\tilde\Phi_\vartheta\dd\rho + (1/k_e)U_{\mathrm{even}} + O(k_e^{-2})$.

\noindent\emph{(d)} $\frac{\mathrm{d}}{\mathrm{d}t}r_\mu(x_m) = \frac{\mathrm{d}}{\mathrm{d}t}f(x_m)
- \int\frac{\mathrm{d}}{\mathrm{d}t}\tilde\Phi(x_m;\vartheta')\dd\mu(\vartheta')$.
\begin{itemize}
  \item Antisymmetry: $K(\vartheta,\vartheta') + K(\vartheta',\vartheta)
    = \frac{1}{T_0}\oint\mathrm{d}\bigl(\tilde\Phi_\vartheta^\T\tilde\Phi_{\vartheta'}\bigr) = 0$.
  \item The area form is Stokes, or Green's theorem for $n=2$.
  \item For $n=1$, every $1$-form on $\R$ is exact, so its loop integrals vanish.
  \item Reversal changes the sign of oriented loop integrals and preserves $\rho$.
\end{itemize}

\noindent\emph{(e)} The second variation of $U_{\mathrm{even}}$ is
\[
  S_{\mathrm{even}}(\vartheta,\vartheta')
  = -\int\Bigl[(\tilde\Phi_{\vartheta'}\cdot\nabla)(r_\mu^\T\tilde\Phi_\vartheta)
    + (r_\mu\cdot\nabla)(\tilde\Phi_{\vartheta'}^\T\tilde\Phi_\vartheta)\Bigr]\dd\rho,
\]
whose second term is symmetric. The second variation of $U_{\mathrm{odd}}$ is $K$. A kernel's antisymmetric part $A$
is a coboundary if and only if $A(\vartheta_0,\vartheta_1) + A(\vartheta_1,\vartheta_2) + A(\vartheta_2,\vartheta_0)
= 0$ for all triples.

The remainder of~(b) is not differentiated in $\mu$. Instead the argument uses a circulation around a small triangle of
measures.
\begin{enumerate}[label=\arabic*.]
  \item Fix $\vartheta_0,\vartheta_1,\vartheta_2$ and $\varepsilon\in(0,1)$. For $(a,b)$ in
    $T_\varepsilon := \{a,b\ge0,\ a+b\le\varepsilon\}$, let
    $\mu_{a,b} := (1-\varepsilon)\mu + (\varepsilon-a-b)\delta_{\vartheta_0} + a\delta_{\vartheta_1}
    + b\delta_{\vartheta_2}$.
  \item The circulation $\oint_{\partial T_\varepsilon}\ip{Q_{\mu_{a,b}}}{\mathrm{d}\mu_{a,b}}$ vanishes when $Q_\mu$
    is a first variation plus a constant $c(\mu)$, because $\mathrm{d}\mu_{a,b}$ has zero mass. In particular it
    vanishes for $-(1/k_e)V_L$.
  \item $U_{\mathrm{even}}+U_{\mathrm{odd}}$ is quadratic in $\mu$, so its circulation equals
    $-(\varepsilon^2/2)\cdot\sum_{\mathrm{cyc}}(A_{\mathrm{even}}+2K)(\vartheta_i,\vartheta_{i+1})$ at the centroid of
    $T_\varepsilon$.
  \item The $O(k_e^{-3})$ remainder of~(b) contributes at most $6\varepsilon\cdot\sup\abs{R} = O(\varepsilon k_e^{-3})$.
  \item Suppose that, for all large $k_e$, $P_\mu$ were a first variation up to $O(k_e^{-3})$, with constants uniform
    in $k_e$. Then $(\varepsilon^2/(2k_e^2))\abs{\mathrm{cyc}}\le C\varepsilon k_e^{-3}$, that is,
    $\abs{\mathrm{cyc}}\le 2C/(\varepsilon k_e)$. The cyclic sum does not depend on $k_e$, so letting
    $k_e\to\infty$ at fixed $\varepsilon$ gives $\mathrm{cyc}=0$ at the centroid measure of the triangle; letting
    $\varepsilon\to0$, since $A_{\mathrm{even}}+2K$ is affine in the measure, gives $\mathrm{cyc}=0$ at $\mu$.
\end{enumerate}
This is checked numerically: for $\varepsilon=0.05$ and $n=2$ the circulation equals the centroid formula to all
printed digits.
\end{proof}

\checkedby{\code{tests/test_averaging.py} in the accompanying repository, using \code{src/core/averaging.py}; seven
tests: four for Proposition H.7, one for Remark H.0, two for Remark H.7a. The checks of Proposition H.7 are listed in
Table~\ref{tab:finitegain:checks}.}

\begin{table}[htbp]
\centering
\small
\caption[Checks of Proposition H.7]{Checks of Proposition H.7 (\code{tests/test_averaging.py}).}
\label{tab:finitegain:checks}
\begin{tabularx}{\textwidth}{@{}>{\raggedright\arraybackslash}p{0.30\textwidth}Y@{}}
\toprule
check & result \\
\midrule
expansion: relative error of $k_e^2(P_\mu - P_0)$ against $U_{\mathrm{even}}+U_{\mathrm{odd}}$
  & halves when $k_e$ doubles. At $k_e = 40\to80$ it goes $0.54\to0.25$ for $n=1$, $0.038\to0.019$ for $n=2$
    counter-clockwise, and $0.022\to0.011$ clockwise. A Richardson step at $k_e=80$ and $160$ brings $n=1$ within
    3\% (asserted) \\
\addlinespace
Stokes: loop-integral $K$ against the area integral over the enclosed ellipse
  & agreement to $7\cdot10^{-15}$ relative (asserted at $10^{-6}$); $K$ antisymmetric; $K\equiv0$ for $n=1$ \\
\addlinespace
parity under orbit reversal (same $\rho$)
  & $U_{\mathrm{even}}$ unchanged and $U_{\mathrm{odd}}$ flipped, to $2\cdot10^{-15}$ relative (asserted at
    $10^{-10}$) \\
\addlinespace
non-reciprocity: the \emph{exact} second variation $S$ ($k_e=80$) against the prediction
  & the full matrix $k_e^2S + k_eS_0$ is within 5\% of $S_{\mathrm{even}}+K$. The cyclic sums of the antisymmetric
    part agree to 3\%: $n=1$ gives $-0.0711$ against $-0.0704$, all from the state shift; $n=2$ gives $-0.0471$
    against $-0.0464$ (state shift $-0.089$, circulation $+0.043$) \\
\bottomrule
\end{tabularx}
\end{table}

\subsection{The oscillation correction and the orientation of the reference orbit}

\begin{rmk}{Remark}{H.7a}[the full first-order correction, and a test in the real closed loop]
Proposition H.7 freezes the measure over a period. In the closed loop the parameters also oscillate within each period,
by $O(1/k_e)$.

\emph{The parameter equation.} In the scaling $\sigma=\sigma_0/k_e$, write
\begin{align}
  \dot\Theta &= \frac{1}{k_e}H(\Theta,t) + \frac{1}{k_e^2}H_2(\Theta,t) + O(k_e^{-3}),\label{eq:finitegain:Theta}\\
  H_i(\Theta,t) &:= G\nabla_\vartheta\bigl[r_\Theta(x_m(t))^\T\tilde\Phi(x_m(t);\vartheta_i)\bigr]
    - \sigma_0(\vartheta_i-\vartheta^0).\notag
\end{align}

\emph{Second-order averaging} (Sanders--Verhulst--Murdock~\cite[\S2.9]{svm2007}) gives the dynamics of the period
mean $Y$:
\begin{equation}\label{eq:finitegain:Y}
  \dot Y = \frac{1}{k_e}\langle H\rangle(Y)
    + \frac{1}{k_e^2}\Bigl[\langle H_2\rangle(Y) + \bigl\langle D_\Theta H(Y,t)\cdot u(Y,t)\bigr\rangle\Bigr]
    + O(k_e^{-3}).
\end{equation}
\begin{itemize}
  \item $u$ is the zero-mean antiderivative in $t$ of $H-\langle H\rangle$.
  \item $\langle H_2\rangle$ is the frozen-measure correction of Proposition H.7.
  \item This is formal. The averaging theorem gives it on finite slow-time horizons.
  \item \code{averaging.oscillation_correction} computes the oscillation term
    $b_{\mathrm{osc}} := \langle D_\Theta H\cdot u\rangle$.
\end{itemize}

Three properties follow.
\begin{itemize}
  \item \textbf{Parity and value.} The oscillation term $b_{\mathrm{osc}}$ is odd under reversal of the orbit
    (\code{test_oscillation_correction_is_odd_under_orbit_reversal}). Its value matches an independent one-period
    near-identity flow at $k_e=80$: the test asserts 3\%, and 0.7\% is measured
    (\code{test_oscillation_correction_matches_a_one_period_flow}).
  \item \textbf{Size.} Relative to the frozen-measure terms it scales like $\abs{G}T_0$, the parameter displacement
    per period times $k_e$. It is negligible when adaptation is slow within a period, which is the small-gain limit in
    the fixed-$k_e$ discussion of Conjecture H.6. It is not negligible as $k_e\to\infty$ at fixed $G$. Both terms are
    $O(k_e^{-2})$ and uniformly bounded, so the order of the route to Conjecture H.6 is unchanged.
  \item \textbf{Non-reciprocity survives.} The only orientation-even part of the full correction is
    $U_{\mathrm{even}}$, and its antisymmetric second variation is not a coboundary (Proposition H.7(e)). Suppose the
    full first-order dynamics were reciprocal for both orientations. Then the average of the two antisymmetric parts,
    $A_{\mathrm{even}}$, would be a coboundary, which it is not. So the full first-order averaged dynamics is
    non-reciprocal for at least one orientation.
\end{itemize}

\emph{A prediction tested in the real closed loop} (\code{exp_orbit_orientation.py} $\to$
\code{results/orbit_orientation/}).
\begin{itemize}
  \item \textbf{The prediction.} Traverse the same elliptic reference both ways. The occupation measure is the same,
    so $\Aavg$, and any functional of $\rho$, cannot distinguish the two runs. The first-order theory predicts that the
    learned kernels differ by
    $\Theta_+ - \Theta_- \approx -2J^{-1}(b_{\mathrm{odd}} + b_{\mathrm{osc}}) = O(1/k_e)$. Here $J$ is the Jacobian
    of the leading averaged drift at the equilibrium.
  \item \textbf{The setting.} $n=2$, $m=6$, $G=\operatorname{diag}(4,4,1,1,1)$, $\sigma=2/k_e$, no dither, and the
    same initial parameters for both runs. The runs converged to machine precision. The six kernels start spread
    out (kernel spread $1.05$) and converge to a common configuration (spread at most $5.6\cdot10^{-17}$ at the
    end of every run), so the learned network is effectively a single kernel: the test probes the orientation
    effect at a collapsed equilibrium of the $m=6$ system.
\end{itemize}

\begin{center}
\small
\captionof{table}[Orientation dependence of the learned kernels]{Orientation dependence of the learned kernels in the
real closed loop (\code{results/orbit_orientation/table_sigma2_slow16.md}). The size ratio is measured over
predicted.}
\label{tab:finitegain:orientation}
\begin{tabular}{@{}ccccccc@{}}
\toprule
 & & \multicolumn{2}{c}{frozen measure only (H.7)} & \multicolumn{2}{c}{full (H.7 + oscillation)} & full \\
\cmidrule(lr){3-4}\cmidrule(lr){5-6}\cmidrule(l){7-7}
$k_e$ & measured $\norm{\Theta_+-\Theta_-}$ & size ratio & cosine & size ratio & cosine & vector error \\
\midrule
10 & $5.7\cdot10^{-3}$ & 2.67 & 0.76 & 0.94 & 0.93  & 0.36 \\
20 & $2.8\cdot10^{-3}$ & 2.41 & 0.83 & 0.97 & 0.98  & 0.20 \\
40 & $1.4\cdot10^{-3}$ & 2.29 & 0.86 & 0.97 & 0.994 & 0.11 \\
\bottomrule
\end{tabular}
\end{center}

\begin{itemize}
  \item The measured difference scales like $k_e^{-1.01}$.
  \item At $k_e=40$ the full first-order theory predicts it to within 3\% in size and 11\% as a vector (cosine 0.994).
  \item The vector error halves with each doubling of $k_e$: 36\%, 20\% and 11\% at $k_e=10$, 20 and 40.
  \item The frozen-measure terms alone account for less than half of it.
  \item A first attempt with a weaker leak ($\sigma=0.5/k_e$) and a shorter horizon had not converged: its residual
    drift was a third of the difference. It is not used.
\end{itemize}

The direction in which the reference orbit is traversed is therefore a property of the excitation that the learned
kernels retain, at order $1/k_e$. The first-order theory predicts its effect.
\end{rmk}
\status{The second-order averaged equation~\eqref{eq:finitegain:Y} is formal; the averaging theorem gives it on
finite slow-time horizons. The oscillation term is checked against an independent one-period near-identity flow
(within 3\% at $k_e=80$) and is odd under orbit reversal (to $10^{-6}$). The full first-order theory predicts the
orientation dependence of the learned kernels in the real closed loop to 3\% in size and 11\% as a vector
(cosine 0.994, $k_e=40$).}
\checkedby{\code{tests/test_averaging.py} (the two tests named above); \code{exp_orbit_orientation.py} $\to$
\code{results/orbit_orientation/table_sigma2_slow16.md}.}

\subsection{Reciprocity at order \texorpdfstring{$k_e^{-2}$}{k\_e\^{}-2}: a pairwise criterion and its
characterisation}

\begin{thm}{Theorem}{I}[reciprocity at order $k_e^{-2}$: a pairwise criterion and its characterisation]
Assume the setting of Proposition H.7: saturated Gaussian features $\tilde\Phi(x;\vartheta)=a\varphi_\xi(x)$ with
$\xi:=(c,s)$, $f\in C^2(\R^n)$ autonomous, $x_m\in C^2$ $T_0$-periodic with oriented orbit $\mathcal{O}$ and
occupation measure $\rho$, and a frozen $\mu\in\Prob(\R^p)$. Let $S$ be the $k_e^{-2}$ coefficient of the second
variation of $P_\mu$, and $A_\mu:=S-S^\T=A_{\mathrm{even}}+2K$ its antisymmetric part, as in H.7(e). Call $\mu$
\emph{reciprocal at order $k_e^{-2}$} if $A_\mu\equiv0$ (the $1$-form that the $k_e^{-2}$ drift defines is closed at
$\mu$). Note that $A_\mu$ is affine in $\mu$, because $r_\mu$ is and $K$ does not depend on $\mu$.
\begin{enumerate}[label=(\roman*)]
  \item \emph{Pairwise criterion.} $A_\mu$ is a coboundary $a(\vartheta)-a(\vartheta')$ if and only if
    $A_\mu\equiv0$. Consequently, by Proposition H.7(e), whose hypothesis concerns a neighbourhood of $\mu$ (its
    triangle of measures), $P_\nu$ can be a first variation of a functional of $\nu$ up to $O(k_e^{-3})$, with
    constants uniform in $k_e$ and in $\nu$ near $\mu$, only if $A_\nu\equiv0$ for every $\nu$ in a neighbourhood of
    $\mu$.
  \item \emph{Structure.} With $Dr$ the Jacobian of $r_\mu$, $(Dr)_{kl}=\partial_lr_k$, and
    $[X,Y]:=(X\cdot\nabla)Y-(Y\cdot\nabla)X$ the Lie bracket of vector fields on $\R^n$,
    \begin{equation}\label{eq:finitegain:AevenLie}
      A_{\mathrm{even}}(\vartheta,\vartheta')=-\int\Bigl\{\tilde\Phi_\vartheta^\T(Dr_\mu-Dr_\mu^\T)\tilde\Phi_{\vartheta'}
        +r_\mu^\T[\tilde\Phi_{\vartheta'},\tilde\Phi_\vartheta]\Bigr\}\dd\rho.
    \end{equation}
    For the Gaussian features $A_\mu$ is a bilinear form in the weights:
    $A_\mu(\vartheta,\vartheta')=a^\T B_\mu(\xi,\xi')a'$ with
    \begin{gather}
      \begin{aligned}
        B_\mu(\xi,\xi'):={}&-\int(Dr_\mu-Dr_\mu^\T)\varphi_\xi\varphi_{\xi'}\dd\rho
          -\int r_\mu(\nabla\varphi_\xi)^\T\varphi_{\xi'}\dd\rho\\
        &+\int(\nabla\varphi_{\xi'})r_\mu^\T\varphi_\xi\dd\rho+2\kappa(\xi,\xi')I_n,
      \end{aligned}\label{eq:finitegain:B}\\
      \kappa(\xi,\xi'):=\Bigl\langle\varphi_\xi(x_m(t))\cdot\frac{\mathrm{d}}{\mathrm{d}t}\varphi_{\xi'}(x_m(t))\Bigr\rangle_t,
        \label{eq:finitegain:kappa}
    \end{gather}
    so $A_\mu\equiv0$ if and only if $B_\mu(\xi,\xi')=0$ for all admissible $\xi,\xi'$. The first term is the curl of
    the residual paired with the features, the next two are the residual paired with the Lie bracket of the feature
    fields, and the last is the circulation kernel $K(\vartheta,\vartheta')=(a^\T a')\kappa(\xi,\xi')$. The
    performative part is bounded by the residual and by the antisymmetric part of its Jacobian along the orbit:
    \[
      \abs{A_{\mathrm{even}}(\vartheta,\vartheta')}\le\abs{a}\abs{a'}\Bigl[\int\norm{Dr_\mu-Dr_\mu^\T}\dd\rho
        +\frac{2e^{-1/2}}{s_{\min}}\int\abs{r_\mu}\dd\rho\Bigr].
    \]
  \item \emph{$n=1$.} $A_\mu\equiv0$ if and only if $r_\mu=0$ $\rho$-a.e.\ on the orbit. So the one-dimensional loop
    is reciprocal at order $k_e^{-2}$, at a frozen $\mu$, exactly when $\mu$ interpolates $f$ along the orbit, and its
    non-reciprocity is bounded by the residual:
    $\abs{A_\mu(\vartheta,\vartheta')}\le(2e^{-1/2}/s_{\min})\abs{a}\abs{a'}\int\abs{r_\mu}\dd\rho$.
  \item \emph{$n\ge2$ at perfect fit.} Suppose $r_\mu=0$ $\rho$-a.e. Then the Lie-bracket terms vanish and
    \[
      B_\mu(\xi,\xi')=-\int(Dr_\mu-Dr_\mu^\T)\varphi_\xi\varphi_{\xi'}\dd\rho+2\kappa(\xi,\xi')I_n,
    \]
    an antisymmetric matrix plus a multiple of the identity. Hence $A_\mu\equiv0$ if and only if both
    \begin{enumerate}[label=(\alph*)]
      \item $K\equiv0$, which holds if and only if the $1$-current of the oriented orbit vanishes,
        $T_{\mathcal{O}}:=\int_0^{T_0}\delta_{x_m(t)}\,\dot x_m(t)\dd t=0$ (an $\R^n$-valued measure;
        $\int v\cdot\dd T_{\mathcal{O}}=\oint_{\mathcal{O}}v\cdot\dd x$); for $n=2$, iff the winding-number function
        $w_{\mathcal{O}}$ of the orbit vanishes a.e.; and
      \item $Dr_\mu=Dr_\mu^\T$ $\rho$-a.e.: the Jacobian of the residual has no antisymmetric part along the orbit.
    \end{enumerate}
    Vanishing of $r_\mu$ on the orbit constrains only its tangential derivative, so (b) is a genuine extra condition:
    for $f:=\fhat_\mu+he_1$ with $h=0$ on the orbit and $\partial_\perp h\neq0$,
    $A_\mu-2K=-\int\varphi_\xi\varphi_{\xi'}\partial_\perp h\dd\rho\cdot(e_1e_\perp^\T-e_\perp e_1^\T)\neq0$. When
    $r_\mu$ vanishes on a neighbourhood of the orbit, $Dr_\mu=0$ there and $A_\mu=2K$: the semi-gradient defect has
    vanished and the loop is non-reciprocal exactly when the orbit carries circulation, through the orbit-area kernel
    alone.
  \item \emph{General $n\ge2$.} In any connected finite-dimensional real-analytic family of measures, for example
    atomic measures with $J$ atoms parametrised by their weights and lifted positions, the set of reciprocal measures
    is either the whole family or a closed real-analytic subvariety of measure zero. The instances of
    \code{tests/test_averaging.py} ($n=1$ and the $n=2$ ellipse) have $A_\mu\not\equiv0$, with cyclic sums $-0.0704$
    and $-0.0464$ of the predicted kernel.
  \item \emph{Nowhere a gradient flow.} For every $n$, every $(f,x_m)$ and every $\mu$, every neighbourhood of $\mu$
    contains measures $\nu$ with $A_\nu\not\equiv0$: along the segments $\nu_\delta:=(1-\delta)\mu+\delta\delta_{\vartheta'}$
    either the zero feature $\vartheta'=\vartheta_z$ or a feature with $a'\neq0$ gives $A_{\nu_\delta}\not\equiv0$ for
    all small $\delta>0$. Hence, by (i), there is no open set of measures on which $P_\nu$ is a first variation up to
    $O(k_e^{-3})$ with uniform constants: the exactly averaged closed loop is nowhere the $\WG$-gradient flow of a
    free energy at order $k_e^{-2}$. The reciprocal measures of (iii) and (iv) are coincidences without interior.
\end{enumerate}
\end{thm}
\status{Proved; every formula checked against the generic second-variation code and, by the round-5 referee, against
the exact second variation; the characterisations checked on perfect-fit, curl and zero-current instances. Novelty
(literature audit): new; the analytic tools (injectivity of Gaussian convolution; Stokes with winding numbers) are
classical.}

\begin{proof}
\emph{(i)} Let $\vartheta_z:=(\alpha,\gamma,\kappa)$ with $\alpha=0$, an admissible point of $\R^p$, so $a=0$ and
$\tilde\Phi(\cdot;\vartheta_z)\equiv0$ as a function of $x$; its $x$-derivatives vanish too. Every term of
$A_{\mathrm{even}}(\vartheta,\vartheta_z)$ contains $\tilde\Phi_{\vartheta_z}$ or $(\tilde\Phi_{\vartheta_z}\cdot\nabla)$,
and $K(\vartheta,\vartheta_z)=\langle\tilde\Phi_\vartheta^\T\frac{\mathrm{d}}{\mathrm{d}t}\tilde\Phi_{\vartheta_z}\rangle=0$;
so $A_\mu(\vartheta,\vartheta_z)=0$ for every $\vartheta$. If $A_\mu(\vartheta,\vartheta')=a(\vartheta)-a(\vartheta')$,
then $a(\vartheta)=a(\vartheta_z)$ for every $\vartheta$, so $A_\mu\equiv0$. The converse is trivial. For the
consequence, apply H.7(e) at any $\nu$ near $\mu$ with $\vartheta_2:=\vartheta_z$ in its triangle of measures: the
cyclic sum over $(\vartheta_0,\vartheta_1,\vartheta_z)$ is $A(\vartheta_0,\vartheta_1)$, evaluated at the centroid
measure of the triangle, which tends to $\nu$ as the triangle's size $\varepsilon\to0$; $A$ is affine in the measure,
so $A_\nu(\vartheta_0,\vartheta_1)=0$.

\noindent\emph{(ii)} Expand
$(\tilde\Phi'\cdot\nabla)(r^\T\tilde\Phi)=\sum_{k,l}\tilde\Phi'_l(\partial_lr_k)\tilde\Phi_k
+\sum_{k,l}\tilde\Phi'_lr_k\partial_l\tilde\Phi_k=\tilde\Phi^\T(Dr)\tilde\Phi'+r^\T(\tilde\Phi'\cdot\nabla)\tilde\Phi$,
and subtract the expression with $\vartheta$ and $\vartheta'$ exchanged. For $\tilde\Phi=a\varphi$,
$(\tilde\Phi'\cdot\nabla)\tilde\Phi=a(a'^\T\nabla\varphi)\varphi'$, so
$r^\T[\tilde\Phi',\tilde\Phi]=(r^\T a)(a'^\T\nabla\varphi)\varphi'-(r^\T a')(a^\T\nabla\varphi')\varphi$, and
$K(\vartheta,\vartheta')=(a^\T a')\kappa(\xi,\xi')$. Collect the bilinear form in $(a,a')$. Since $a$ ranges over the
open cube $(-a_{\max},a_{\max})^n$, $A_\mu\equiv0$ iff $B_\mu\equiv0$. The bound uses $\varphi\le1$ and
$\abs{\nabla\varphi}\le e^{-1/2}/s_{\min}$.

\noindent\emph{(iii)} For $n=1$, $Dr-Dr^\T=0$, and for equal widths $s$,
$\varphi_{c'}D\varphi_c-\varphi_cD\varphi_{c'}=\varphi_c\varphi_{c'}(c-c')/s^2$, so
$B_\mu(\xi,\xi')=-((c-c')/s^2)\int r_\mu\varphi_c\varphi_{c'}\dd\rho$. Now
$\varphi_c\varphi_{c'}=e^{-(c-c')^2/(4s^2)}\cdot e^{-(x-m)^2/s^2}$ with $m:=(c+c')/2$. Put $\nu:=r_\mu\rho$, a finite
signed measure with compact support. The function $m\mapsto\int e^{-(x-m)^2/s^2}\dd\nu(x)$ extends to an entire
function. If $B_\mu\equiv0$ on admissible pairs, it vanishes for all $m\in(-c_{\max},c_{\max})$ (take $c=m-\delta$,
$c'=m+\delta$ with $0<\delta$ small), hence everywhere; its Fourier transform is the product of the Gaussian's, which
has no zeros, with $\hat\nu$, so $\nu=0$, that is $r_\mu=0$ $\rho$-a.e. Conversely, if $\nu=0$ then
$A_{\mathrm{even}}(\vartheta,\vartheta')=-\int(\varphi'D\varphi-\varphi D\varphi')aa'\dd\nu=0$ for all widths.

\noindent\emph{(iv)} If $r_\mu=0$ $\rho$-a.e.\ the two Lie-bracket integrals vanish, which gives the displayed
$B_\mu$; the first term is an antisymmetric matrix and the second a multiple of $I_n$, so $B_\mu\equiv0$ iff both
vanish (take the trace, which is $2n\kappa$, and the antisymmetric part). The first vanishes identically iff the
matrix-valued measure $(Dr_\mu-Dr_\mu^\T)\rho$ is zero, by the argument of (iii) applied entrywise with equal widths
(the entries of $\int(Dr-Dr^\T)\varphi_c\varphi_{c'}\dd\rho$ are Gaussian convolutions of the entries of that
measure). For (a): with equal widths, $\varphi_\xi\varphi_{\xi'}=e^{-\abs{c-c'}^2/(4s^2)}e^{-\abs{x-m}^2/s^2}$ and
$\nabla\varphi_{\xi'}=-((x-c')/s^2)\varphi_{\xi'}$, so
\[
  \kappa(\xi,\xi')=-\frac{1}{s^2}e^{-\abs{c-c'}^2/(4s^2)}
  \bigl\langle e^{-\abs{x_m-m}^2/s^2}(x_m-c')\cdot\dot x_m\bigr\rangle_t.
\]
Write $x_m-c'=(x_m-m)+(c-c')/2$. The first part integrates to zero, being the time derivative of
$-(s^2/2)e^{-\abs{x_m-m}^2/s^2}$ along a period. Hence
\begin{gather*}
  \kappa(\xi,\xi')=-\frac{1}{2s^2}e^{-\abs{c-c'}^2/(4s^2)}(c-c')\cdot F_s(m),\\
  F_s(m):=\bigl\langle e^{-\abs{x_m(t)-m}^2/s^2}\dot x_m(t)\bigr\rangle_t
  =\frac{1}{T_0}\int e^{-\abs{x-m}^2/s^2}\dd T_{\mathcal{O}}(x),
\end{gather*}
the Gaussian convolution of the vector measure $T_{\mathcal{O}}$. If $\kappa\equiv0$ on admissible pairs, then
$F_s(m)\cdot v=0$ for every $m$ with $\norm{m}<c_{\max}$ and every direction $v$ (take $c=m-\delta v$,
$c'=m+\delta v$), so $F_s\equiv0$ on that ball; $F_s$ is entire, so $F_s\equiv0$ on $\R^n$, and the Fourier argument
gives $T_{\mathcal{O}}=0$. Conversely,
$\kappa(\xi,\xi')=\frac{1}{T_0}\int\varphi_\xi\nabla\varphi_{\xi'}\cdot\dd T_{\mathcal{O}}=0$ for all widths when
$T_{\mathcal{O}}=0$. For $n=2$, $\oint_{\mathcal{O}}v\cdot\dd x=\int w_{\mathcal{O}}\operatorname{curl}v\dd A$ for
every smooth $v$ (Stokes with winding numbers), and every smooth compactly supported $g$ is $\operatorname{curl}v$ for
some smooth $v$, so $T_{\mathcal{O}}=0$ iff $w_{\mathcal{O}}=0$ a.e.\ ($T_{\mathcal{O}}$ is compactly supported, so
$v$ need not be). For the example, $r=he_1$ with $h=0$ on $\mathcal{O}$ gives $(Dr-Dr^\T)_{1l}=\partial_lh$ and the
only non-zero entries are the normal derivative; the Lie-bracket terms vanish because $r=0$ on $\mathcal{O}$.

\noindent\emph{(v)} The saturated features are real-analytic in the lifted coordinates and extend holomorphically to
a complex neighbourhood of $\R^p$ ($\tanh$ and the logistic function are holomorphic on the strips
$\abs{\operatorname{Im}}<\pi a_{\max}/2$, $\pi c_{\max}/2$, $\pi$, where $s$ stays away from $0$), uniformly for $x$
on the compact orbit; integrals over the orbit of such integrands are real-analytic in the parameters, and $r_\mu$ is
affine in the weights. The zero set of a real-analytic function on a connected real-analytic manifold is either the
whole manifold or a closed set of Lebesgue measure zero. The instances are those of\newline
\code{test_exactly_averaged_dynamics_is_non_reciprocal_for_n1_and_n2}.

\noindent\emph{(vi)} Since $r_\nu$ is affine in $\nu$, $B_{\nu_\delta}=B_\mu+\delta B^{\mathrm{even}}[g]$ with
$g:=\fhat_\mu-\tilde\Phi_{\vartheta'}$, where $B^{\mathrm{even}}[g]$ denotes the three residual terms of $B$ with $g$
in place of $r_\mu$. The trace of $B^{\mathrm{even}}[g]$ at equal widths is
$-(1/s^2)(c-c')\cdot\int g\varphi_c\varphi_{c'}\dd\rho$ (the curl term is traceless). If $A_{\nu_\delta}\equiv0$ for
all small $\delta$, then $B^{\mathrm{even}}[g]\equiv0$, so $(c-c')\cdot F[g\rho]((c+c')/2)=0$ for all admissible
$c,c'$, where $F[g\rho]$ is the Gaussian convolution of the vector measure $g\rho$; as in (iv), $g\rho=0$, that is
$g=0$ $\rho$-a.e. For $\vartheta'=\vartheta_z$ this says $\fhat_\mu=0$ on the orbit; for $\vartheta'$ with $a'\neq0$
it then says $\tilde\Phi_{\vartheta'}=0$ on the orbit, which is false since a Gaussian has no zeros. So one of the two
segments has $A_{\nu_\delta}\not\equiv0$ for all small $\delta>0$.
\end{proof}
\checkedby{\code{tests/test_nonreciprocity.py}, using \code{src/core/nonreciprocity.py} against
\code{src/core/averaging.py}: the bilinear form $a^\T B_\mu a'$ reproduces the antisymmetric part of the generic
second-variation coefficient ($10^{-10}$, $n=1$ and $2$); $A_\mu(\vartheta,\vartheta_z)=0$ ($10^{-14}$) and the cyclic
sum over $(\vartheta_0,\vartheta_1,\vartheta_z)$ equals $A_\mu(\vartheta_0,\vartheta_1)$; for $n=1$ with
$f:=\fhat_\mu$ (perfect fit) $A_\mu\equiv0$ (below $10^{-13}$) while a non-zero residual gives $A_\mu\neq0$, equal to
the closed form; the closed form of $\kappa$ ($10^{-10}$), the second component of $F_s(m)$ at $m=(0.5,0)$, $s=0.7$ is
positive on the ellipse, and $\kappa\equiv0$ on a segment traversed back and forth (below $10^{-12}$); with
$f\equiv\fhat_\mu$ on $\R^2$ (so $Dr_\mu\equiv0$), $B_\mu=2\kappa I$ ($10^{-13}$), $A_\mu\neq0$ on the ellipse and
$A_\mu\equiv0$ on the segment; with $f:=\fhat_\mu+he_1$, $h=(x_1/1.2)^2+(x_2/0.8)^2-1$ (zero on the ellipse, non-zero
normal derivative), $B_\mu-2\kappa I$ equals $-\int\varphi\varphi'\partial_2h\dd\rho\cdot(e_1e_2^\T-e_2e_1^\T)$
($10^{-12}$) and is not small; the positivity witness $f\equiv0$, $\mu=\delta_{\bar\vartheta}$, for which
$A_\mu(\vartheta_0,\vartheta_1)=a_0a_1\bar a\,((c_0-c_1)/s^2)\int\varphi_{\bar c}\varphi_{c_0}\varphi_{c_1}\dd\rho$ has
a strictly positive integral. The round-5 referee re-derived every formula and checked $a^\T B_\mu a'$ against the
exact second variation of the averaged potential (Richardson in $k_e$) to $2\cdot10^{-4}$--$9\cdot10^{-4}$ relative,
including parameters at the saturation limits.}

\begin{rmk}{Remark}{I.1}[what Theorem I changes in the reading of Proposition H.7]\mbox{}
\begin{itemize}
  \item The non-reciprocity of H.7(e) is no longer an instance-by-instance check. Pointwise, for $n=1$ it holds at
    every frozen $\mu$ that does not interpolate $f$ along the orbit; for $n\ge2$ at perfect fit it holds unless the
    orbit carries no circulation and the residual's Jacobian is symmetric along the orbit; in general it fails only on
    a null set of any connected analytic family that contains one non-reciprocal member. Globally, (vi) shows that no
    open set of measures is reciprocal, for any $n$: the averaged loop is nowhere a gradient flow at order $k_e^{-2}$.
  \item The two defects scale differently. The performative (semi-gradient) defect is bounded by the residual and the
    antisymmetric part of its Jacobian along the orbit, so it is small at a well-fitting equilibrium whose residual is
    smooth across the orbit; the circulation defect does not depend on the fit at all.
  \item For the $n=1$ Langevin experiments and Conjecture H.6: the non-gradient correction that falsifies H.6 at
    fixed $k_e$ is, by (iii), bounded by the residual, which at the regularised optimum $\mu^*_\tau$ is the
    regularisation bias. Both the leading drift and the correction are proportional to the residual there, so the
    correction is $O(1/k_e)$ relative to the leading drift and $O(k_e^{-2})$ absolute, as in Remark H.0.
  \item The pairwise criterion also simplifies the test of reciprocity: a single pair $(\vartheta,\vartheta')$ with
    $B_\mu(\xi,\xi')\neq0$ suffices.
\end{itemize}
\end{rmk}

\subsection{Input-to-state stability of \texorpdfstring{$\Aavg$}{(A)} and the route to Conjecture H.6}

\begin{thm}{Proposition}{H.8}[input-to-state stability of $\Aavg$]
Assume $\sigma,\tau>0$ and the setting of Theorem H. Let $\mu_t$ solve the perturbed equation
\begin{equation}\label{eq:finitegain:perturbed}
  \partial_t\mu = \nabla\cdot\bigl(\mu(G\nabla V[\mu] - b_t)\bigr) + \tau\nabla\cdot(G\nabla\mu)
\end{equation}
from $\mu_0\in\Prob_2$ with $F_\tau(\mu_0)<\infty$. Here $b_t$ is a velocity perturbation with
$\beta(t)^2 := \int\norm{b_t}^2_{G^{-1}}\dd\mu_t\le\bar\beta^2$. Assume the chain rule of~\cite[Lemma~2.2]{S3}
holds along this solution:
\[
  \frac{\mathrm{d}}{\mathrm{d}t}F_\tau(\mu_t) = \int\ip{\nabla(\delta F_\tau/\delta\mu)(\mu_t)}{v_t}\dd\mu_t,
\]
in whitened coordinates, with $v_t$ the velocity. It is expected for smooth bounded perturbations (parabolic regularity and the
entropy-dissipation estimate), but it is not verified here. Then
\begin{align}
  \frac{\mathrm{d}}{\mathrm{d}t}\bigl(F_\tau(\mu_t)-F^*\bigr)
    &\le -\tau\rho_\tau\bigl(F_\tau(\mu_t)-F^*\bigr) + \tfrac12\beta(t)^2,\label{eq:finitegain:iss1}\\
  F_\tau(\mu_t)-F^* &\le e^{-\tau\rho_\tau t}\bigl(F_\tau(\mu_0)-F^*\bigr) + \frac{\bar\beta^2}{2\tau\rho_\tau},
    \label{eq:finitegain:iss2}\\
  \limsup_{t\to\infty}\WG(\mu_t,\mu^*_\tau) &\le \frac{\bar\beta}{\tau\rho_\tau}.\label{eq:finitegain:iss3}
\end{align}
\end{thm}
\status{Proved, assuming the chain rule for the perturbed flow. A proof only; not checked in the repository.
Novelty (literature audit): elementary given~\cite{S3}.}

\begin{proof}
Work in $\zeta=G^{-1/2}\vartheta$, and let
$\xi_t := \nabla(\delta F_\tau/\delta\mu)(\mu_t) = \nabla\tilde V[\mu_t] + \tau\nabla\log\mu_t$. The velocity is
$-\xi_t+\tilde b_t$, with $\norm{\tilde b_t} = \norm{b_t}_{G^{-1}}$.
\begin{enumerate}[label=\arabic*.]
  \item By the chain rule and Young's inequality,
    \[
      \frac{\mathrm{d}}{\mathrm{d}t}F_\tau(\mu_t) = -\int\norm{\xi_t}^2\dd\mu_t + \int\ip{\xi_t}{\tilde b_t}\dd\mu_t
      \le -\frac12\int\norm{\xi_t}^2\dd\mu_t + \frac12\beta^2.
    \]
  \item With $\hat\mu_t\propto e^{-\tilde V[\mu_t]/\tau}$,
    \[
      \int\norm{\xi_t}^2\dd\mu_t = \tau^2 I(\mu_t|\hat\mu_t) \ge 2\rho_\tau\tau^2 H(\mu_t|\hat\mu_t)
      \ge 2\rho_\tau\tau\bigl(F_\tau(\mu_t)-F^*\bigr),
    \]
    by the LSI of Theorem H(ii) and the entropy sandwich~\cite[Lemma~3.4]{S3}.
  \item Grönwall gives the first two lines.
  \item For the third, $F_\tau(\mu)-F^*\ge\tau H(\mu|\mu^*_\tau)$~\cite[Lemma~3.4]{S3}. $\mu^*_\tau$ satisfies
    $\mathrm{LSI}(\rho_\tau)$, hence Talagrand's inequality $W^2\le(2/\rho_\tau)H(\cdot|\mu^*_\tau)$
    (Otto--Villani~\cite{ov2000}). So $W^2\le(2/(\tau\rho_\tau))(F_\tau-F^*)$, and the limsup follows from the second
    line.\qedhere
\end{enumerate}
\end{proof}

The gain $1/(\tau\rho_\tau)\le e^{\alpha^*/\tau}/\sigma$ is exponentially large in $1/\tau$, like the inverse of the
rate.

\paragraph{Route to Conjecture H.6.} In the scaling of Conjecture H.6:
\begin{itemize}
  \item $\tau\rho_\tau\ge(\sigma_0/k_e)\cdot e^{-2K_\Phi(F_\rho+K_\Phi)/\tau_0}$, so $1/(\tau\rho_\tau) = O(k_e)$.
    \begin{itemize}
      \item In this scaling the whitened proximal Gibbs measures do not depend on $k_e$, so
        $\tau\rho_\tau=\Theta(1/k_e)$.
      \item At the experiment's $\tau_0=0.04$ the constant is $e^{1049}$.
    \end{itemize}
  \item Define $b := (1/k_e^2)\bigl[G\nabla(U_{\mathrm{even}}+U_{\mathrm{odd}}) + b_{\mathrm{osc}}\bigr] + O(k_e^{-3})$.
    The first part is the frozen-measure correction of Proposition H.7; the second is the oscillation term of Remark
    H.7a. Both are bounded uniformly in $\mu$ and $\vartheta$, so $\bar\beta = O(k_e^{-2})$.
  \item Proposition H.8 then confines every solution of $\Aavg+b$ to an asymptotic $W$-ball around $\mu^*_\tau$ of
    radius $O(1/k_e)\cdot e^{2K_\Phi(F_\rho+K_\Phi)/\tau_0}/\sigma_0$. This needs $F_\tau(\mu_0)<\infty$, or a
    $\Prob_2$ start, after the Girsanov regularisation of Theorem H(ii), which applies because $b$ is bounded.
\end{itemize}

Three things are still missing for Conjecture H.6:
\begin{enumerate}[label=\arabic*.]
  \item control of the frozen-measure approximation along the moving $\mu_t$, that is, a near-identity averaging
    transformation for the stochastic McKean--Vlasov loop, uniformly in time;
  \item the mean-field limit of $(\mathrm{L}_\tau)$, since Theorem D covers only (L);
  \item the chain rule that Proposition H.8 assumes for the perturbed flow.
\end{enumerate}

\subsection{Evidence from the real closed loop}

The evidence comes from the real closed loop with $m=32$ and $n=1$, not from the averaged system.
\begin{itemize}
  \item \textbf{Dither is necessary.} Without dither the collapsed start is trapped forever. Its loss stays at 0.166,
    against 0.007 from the spread start (\code{results/langevin_tradeoff}, 1500\,s).
  \item \textbf{Equalisation at moderate temperature.} At $\sigma=0.02$ and $\tau=0.004$ (3 seeds,
    \code{results/langevin_global/table_sigma0.02_tau0.004.md}):
    \begin{itemize}
      \item collapsed and spread starts reach the same loss (ratio 1.03) and tracking RMS (0.0465 against 0.0452);
      \item their final 32-particle configurations, driven by the same noise realisation, are 4.2$\times$ closer in
        $W_2$ over the physical coordinates $(a,c,s)$ than without dither (0.40 against 1.69);
      \item a numerical check shows that two independent 32-point samples of $\mu^*_\tau$ are typically about 0.6
        apart, so 0.40 is at the sampling floor;
      \item the common loss of about 0.11 matches the grid value $\bar L(\mu^*_\tau)=0.11118$. It is dominated by the
        regularisation bias.
    \end{itemize}
  \item \textbf{Metastability at low temperature.} At $\sigma=0.002$ and $\tau=5\cdot10^{-4}$ (table for these
    values in \code{results/langevin_global/}), dither frees the collapsed start ($0.166\to0.048$). But the starts have
    not equalised within 900\,s (loss ratio 5.2). The spread start's $9.2\cdot10^{-3}$ is close to the regularised
    optimum $\bar L(\mu^*_\tau)=8.5\cdot10^{-3}$.
  \item \textbf{Limits of the evidence.}
    \begin{itemize}
      \item The Langevin runs have $n=1$. By Proposition H.7, the averaged dynamics is already non-reciprocal for
        $n=1$, through the state shift.
      \item The only $n\ge2$ closed-loop runs are the deterministic orientation runs of Remark H.7a. They confirm the
        first-order theory quantitatively. No $n\ge2$ run with dither has been made.
      \item The runs use fixed $k_e$, not the scaling of Conjecture H.6.
      \item All Langevin runs use $k_e=10$, below the Tikhonov scale of Remark H.0 for this network
        ($(F_\rho+K_\Phi)/s_{\min}\approx35$). The measured frozen-$\mu$ drift error at $k_e=10$ is 4.2\% ($n=1$), so
        the condition is conservative, but the agreement with $\Aavg$ is empirical.
      \item So the evidence supports Conjecture H.6 qualitatively, but the equalisation time can exceed practical
        horizons at low $\tau$.
    \end{itemize}
\end{itemize}
